\section{Hodge classes on mixed products}\label{sec:mixed}

Fix projective K3 surfaces $S_1,\ldots,S_m$, and put
\[
 V_i=T(S_i)(1),\qquad W_i=C^+(T(S_i)).
\]
Here $V_i$ has weight zero, whereas $W_i$ carries the weight-one
Kuga--Satake structure.  Theorem~\ref{thm:exact-ks} supplies the exact
Clifford correspondences, and Corollary~\ref{cor:self-powers} supplies all
Hodge classes on each individual self-power.  We must still treat tensors
involving several different $V_i$.

The individual assertions determine the projections of the joint Hodge
group, but not the relations between those projections.  We first use
weight-one Hodge morphisms between the Kuga--Satake varieties to compare
the actions of shared simple factors projectively.  The CM theorem then
controls the torus action.  After removing these actions, only certain
central signs remain, arising from the two half-spin representations in
dimension four.  We encode their compatibility by a finite graph and use
the degree-four tensors of Proposition~\ref{prop:auxiliary-spin} to prove
the required relations.  All auxiliary abelian varieties will be chosen
before choosing the tensor automorphism to which this comparison is
applied.

\subsection{Tensor groups and the individual factors}

For a finite collection of the surfaces and abelian varieties used here,
let $G$ be its Hodge group, with Tate objects fixed.  Let $C$ be the group
of tensor automorphisms of its cohomology realizations that preserve
algebraic correspondences, again fixing Tate objects.  Then $G\subset C$.
We use these groups over $\Q$ and extend their coefficients to $\C$ when
decomposing representations.

We recall why $C$ has the required invariant interpretation.  The
Lefschetz standard conjecture holds for surfaces and abelian varieties
and is stable under products; it gives algebraic K\"unneth and Lefschetz
operators in the category under consideration
\cite{KleimanStandard}.  Hodge--Riemann then identifies numerical and
homological equivalence there.  Numerical semisimplicity
\cite{JannsenMotives}, with the symmetry modified to give the ordinary
flip on cohomology, gives a neutral Tannakian category of pure motives
with its cohomology fiber functor \cite{DeligneMilneTannakian}.
Its tensor automorphism group is the group just described.

In weight zero, the $C$-invariants are precisely algebraic tensors.
Indeed, before fixing Tate objects, invariants are morphisms from the
unit motive.  The central weight torus acts by cohomological weight and
maps onto the group acting on Tate objects up to isogeny.  It acts
trivially in weight zero, so fixing Tate objects does not change those
invariants.  Consequently it will suffice to prove
\begin{equation}\label{eq:mixed-group-goal}
 C\big|_{\bigoplus_i V_i}=G\big|_{\bigoplus_i V_i}
\end{equation}
for one finite collection containing the original surfaces and the
auxiliary abelian varieties constructed below.

\begin{lemma}[Individual actions]\label{lem:mixed-individual}
For any such collection and any $c\in C(\C)$, the action of $c$ on $V_i$
belongs to $\Hg(V_i)(\C)$.  Its simultaneous action on $V_i,W_i$ is
induced by a lift to the spin group of the quadratic space $V_i$.
\end{lemma}

\begin{proof}
Corollary~\ref{cor:self-powers}, equivalently the exact-tensor criterion
of \cite[Theorem~1.1]{AKS} applied after
Theorem~\ref{thm:exact-ks}, makes every Hodge tensor on powers of $V_i$
algebraic.  Algebraic degree projectors and the quadratic pairing allow
the usual tensor criterion for the Hodge group.  Hence $c|_{V_i}$ lies
in $\Hg(V_i)$.

Choose a spin lift of this orthogonal action.  The operator it induces
on $W_i$ preserves the exact KS tensor, all right even Clifford
multiplications, and an invariant polarization.  These tensors are
also preserved by $c$: the KS tensor is algebraic, while the right
multiplications are weight-one Hodge endomorphisms of an abelian
variety and therefore algebraic.  After removing the chosen spin lift,
the remaining operator on $W_i$ commutes with all right multiplications,
so it is left multiplication by an even Clifford element $a$.
Preservation of the KS map
\[
 v\longmapsto L_vR_w,\qquad w^2\ne0,
\]
then says that $a$ commutes with every $v\in V_i$.  The even part of
the center of the full Clifford algebra consists of scalars.  Thus the
remaining operator is scalar, and preservation of the polarization
makes it $1$ or $-1$.  Both scalars are supplied by the spin cover.
\end{proof}

Let
\[
 G_0=\Hg\left(\bigoplus_i V_i\oplus\bigoplus_i W_i\right).
\]
The projection from $G_0$ to the Hodge group of the $V_i$ has finite
kernel: the Hodge circle acts on each $W_i$ through the spin lift of
its action on $V_i$.  Thus the $W_i$ do not add a Lie algebra factor.

For each $i$, write $E_i=\End_{\rm Hdg}(V_i)$.  Zarhin's theorem says
that $E_i$ is either totally real or CM \cite{Zarhin}.  We use the
following notation:
\begin{itemize}
\item in the orthogonal case, $F_i=E_i$ and
  $d_i=\dim_{E_i}V_i\ge3$;
\item in the unitary case, $F_i=E_i^+$ is the maximal totally real
  subfield and $d_i=\dim_{E_i}V_i\ge1$.
\end{itemize}
Over $\C$, the corresponding groups and representations are
\begin{equation}\label{eq:mixed-single-groups}
 \begin{array}{c|c|c}
 &\Hg(V_i)_{\C}&(V_i)_{\C}\\ \hline
 \text{orthogonal}&\displaystyle\prod_{p:F_i\hookrightarrow\C}
       \mathrm{SO}_{d_i}
   &\text{one standard space at each }p\\[3pt]
 \text{unitary}&\displaystyle\prod_{p:F_i\hookrightarrow\C}
       \mathrm{GL}_{d_i}
   &\text{a standard space and its dual at each }p .
 \end{array}
\end{equation}
The quadratic form in the orthogonal case, and the underlying quadratic
form of the Hermitian form in the unitary case, has positive index two
at the distinguished real embedding and is negative definite at the
other real embeddings.  We call the distinguished embedding the
\emph{active place}.  The infinitesimal Hodge circle has weights
$2,-2$ once each on $V_i$, and weight zero elsewhere.

\subsection{Shared simple ideals and spin constituents}

The derived Lie algebra of $G_0$ is a subdirect product of the simple
ideals occurring in \eqref{eq:mixed-single-groups}.  It is therefore
a product of simple ideals, each mapping diagonally onto some of the
factor ideals.  We call two factor ideals \emph{shared} if they are
images of the same ideal of $\operatorname{Lie}(G_0)^{\rm der}$.

\begin{lemma}[Possible identifications]\label{lem:mixed-shared}
Shared simple ideals have one of the following forms:
\begin{enumerate}
\item orthogonal factors of the same dimension, with their standard
  modules matched;
\item unitary factors of the same dimension, with their standard
  modules matched up to duality;
\item ideals isomorphic to $\mathfrak{sl}_2$, shared among orthogonal
  factors of dimensions three or four and unitary factors of
  dimension two.
\end{enumerate}
Whenever sharing occurs, the fields $F_i$, viewed in $\R$ through
their active embeddings, agree.  The associated place identifications
are then the identity on that common field.  In the unitary case with
$d_i\ge3$, matching the standard modules as representations of the
shared simple ideal also matches their characters on the connected
center of $G_0$.
\end{lemma}

\begin{proof}
An identification of shared ideals must identify their infinitesimal
Hodge-circle elements.  Their eigenvalues first restrict the
possible representations; the Galois translates of the nonzero
components will then identify their fields and places.
First transport a shared ideal by Galois to an active place.  In an
orthogonal standard module the circle element rotates one plane.  For
$\mathfrak{so}_4=\mathfrak{sl}_2\oplus\mathfrak{sl}_2$, both simple
ideals are active, with weights $1,-1$ on their standard spin modules.
The spin standard module for $\mathfrak{so}_3$ has the same weights.
In an active unitary standard module there is exactly one nonzero
weight, equal to $2$ or $-2$; the simple ideal sees its traceless part.

These elements distinguish the possible identifications of the simple
Lie algebras.  In the accidental isomorphism
$\mathfrak{so}_6\simeq\mathfrak{sl}_4$, the orthogonal half-spin
module has weights $1,1,-1,-1$, whereas a unitary standard module
has the traceless pattern of a single nonzero weight.  They cannot
match.  In type $D_4$, a triality automorphism carrying the standard
module to a half-spin module is likewise excluded: the latter has
only weights $1,-1$, unlike the orthogonal standard module.  The
remaining coincidences are precisely those in the statement.

For completeness, the same argument identifies the fields and the
place labels.  Work with simple ideals over $\overline\Q$, testing
activity after the fixed embedding into $\C$.  For an ideal at a
place $p$ of $F_i$, the Galois elements carrying it to an active
ideal are exactly those carrying $p$ to the active place of $F_i$.
This holds also in dimension four, since both ideals there are active.
The sets agree for shared ideals.  After transporting one shared ideal
to the active place, equality of these sets is equality of the
stabilizers of the active field embeddings.  Thus the embedded fields
agree, and all place matchings are the matching embeddings of that
field.  Sharing is transported to every such place; for an orthogonal
dimension-four block this may initially involve only one of its two
simple ideals.

Finally consider a shared unitary ideal of dimension $d\ge3$ and
choose the standard or its dual so that the simple modules agree.
The exceptional eigenvalue in the traceless circle element determines
the sign of its scalar part.  Hence the central characters of the two
matched modules have the same differential on the central component
of the circle.  This is true after every Galois transport.  A Galois
stable character condition that vanishes on the circle vanishes on
the rational torus generated by its projection.  Applying this to
the difference of the two characters proves equality on the connected
center.  One may pass to powers of characters to descend to the
central torus quotient; its character lattice is torsion free, so
equality of these powers implies equality of the original characters.
For $d=2$, the traceless pattern does not distinguish the sign, which
is why that case remains in the third alternative.
\end{proof}

We will use the constituents of $W_i$ as representations of the joint
group.  Restriction of the full even Clifford representation to an
orthogonal sum gives all choices of spin constituents.  In the
orthogonal case, these are tensor products over the places of $F_i$:
there is one spin choice for an odd-dimensional block and either
half-spin choice for an even-dimensional block.  Multiplicities will
not matter.  In dimension three the choice is the standard module of
one $\mathrm{SL}_2$; in dimension four the two choices are the
standard modules of the two $\mathrm{SL}_2$ factors.

In the unitary case the choices at a place are
\begin{equation}\label{eq:mixed-unitary-spin}
 \bigwedge^r\mathrm{Std}\otimes(\det)^{-1/2},
       \qquad 0\le r\le d_i,
\end{equation}
with independent choices at the different places.  The half-determinant
notation denotes a character on the relevant spin cover.  When $d_i=2$,
these choices are the standard module of $\mathrm{SL}_2$, with zero
connected central weight, and two opposite torus characters.  These
descriptions follow from the exterior-algebra model of spin and the
spin decomposition of an orthogonal sum.

\begin{lemma}[Weight-one comparison]\label{lem:mixed-constituents}
If a simple ideal is shared between factors $i$ and $j$, then $W_i$
and $W_j$ contain isomorphic representations of the joint Hodge group
on which that ideal acts nontrivially.  The required intertwiners and
the projections onto individual copies of these representations are
algebraic correspondences with complex coefficients.
\end{lemma}

\begin{proof}
For a shared $\mathfrak{sl}_2$, choose its standard module at the
given place and a transported shared standard module at every other
place.  Lemma~\ref{lem:mixed-shared} provides the common field and
these transported identifications.  In the other orthogonal types,
choose matching spin constituents at all places.  In the unitary
types, choose matching constituents in
\eqref{eq:mixed-unitary-spin}; their central characters agree by
Lemma~\ref{lem:mixed-shared}.  Thus the tensor constituents have the
same joint Lie algebra representation.  The Hodge group is connected,
so they have the same group representation.

Morphisms between weight-one Hodge structures of complex abelian
varieties are induced by rational homomorphisms of abelian varieties.
After extending coefficients, every intertwiner just used is therefore
algebraic.  The same applies to the weight-one idempotents selecting
the constituents and individual multiplicity copies.
\end{proof}

\subsection{Removing projective and torus actions}

We prove the reduction before constructing the auxiliary varieties.
For this subsection, allow any finite collection containing the original
surfaces and their $W_i$ whose Hodge group $G$ projects onto $G_0$ with
finite kernel.  Write $C$ for its algebraic tensor group as above, and
identify the Lie algebras of $G$ and $G_0$.  The auxiliary varieties
constructed below will satisfy this hypothesis.  Thus the reduction
will apply to the final collection once all of its objects have been
fixed.

\begin{lemma}[Projective compatibility]\label{lem:mixed-projective}
For $c\in C(\C)$, its actions at all occurrences of a shared simple
ideal agree in the corresponding adjoint group.
\end{lemma}

\begin{proof}
By Lemma~\ref{lem:mixed-individual}, $c$ acts on each $W_i$ by a
spin lift of its action on $V_i$.  This has the tensor-product
description of its block spin constituents, up to an overall sign.
Lemma~\ref{lem:mixed-constituents} provides matching constituents
on which a specified shared ideal acts nontrivially.  Their
intertwiner is algebraic, so $c$ has the same action in these two
realizations.  Equality of invertible tensor-product operators
forces equality of their individual factors up to scalars.  A
nontrivial representation of a simple Lie algebra is projectively
faithful on its adjoint group.  Hence the adjoint actions at the
shared ideal agree.
\end{proof}

Let $W_{\rm CM}$ be the sum, in all the weight-one realizations of
the collection under consideration, of the constituents on which
$\operatorname{Lie}(G)^{\rm der}$ acts trivially.  This is a
rational sub-Hodge structure.  It is the realization of a CM abelian
factor, and its projector is algebraic.  The CM Hodge conjecture
\cite[Theorem~1.1]{CM}, applied jointly to these factors and their
powers, therefore gives
\begin{equation}\label{eq:mixed-cm-image}
 C\big|_{W_{\rm CM}}=G\big|_{W_{\rm CM}}.
\end{equation}

\begin{lemma}[Reduction to central actions]\label{lem:mixed-central}
After multiplying $c\in C(\C)$ by an element of $G(\C)$, one may
assume that its compatible adjoint actions through the original
factors are trivial at every simple ideal and that it acts as the
identity on $W_{\rm CM}$.  Under these assumptions all
actions on the $V_i$ outside the dimension-four orthogonal blocks
and the dimension-two unitary blocks can be removed by centers of
the shared simply connected derived factors.
\end{lemma}

\begin{proof}
Lemma~\ref{lem:mixed-projective} gives a single prescribed action
in the product of the adjoint groups, and
\eqref{eq:mixed-cm-image} gives the prescribed action on the CM
part.  The derived group covers the first product, and the connected
center covers the torus image on $W_{\rm CM}$.  These choices can
be made independently.  Removing their product proves the first
assertion.  Continue to write $c$ for the resulting element.

On an odd-dimensional orthogonal block with trivial adjoint action,
the standard action is already the identity.  On an even-dimensional
orthogonal block it is a sign.  In a spin lift this sign is the ratio
of the scalar actions on the two half-spin choices.  For simple
types other than $\mathfrak{sl}_2$, these ratios agree at shared
ideals: apply the weight-one intertwiners to matching tensor words
with either half-spin choice at the specified place.  Interchanging
the two choices in one description only inverts the ratio and does
not change a sign.  The center of the shared simply connected
factor supplies that standard-module sign, so it can be removed.

On a unitary block at $(i,p)$, choose the standard or its dual so
that the simple modules match on each shared ideal, as in
Lemma~\ref{lem:mixed-shared}.  Write its scalar action as $r_{i,p}$;
the opposite choice has action $r_{i,p}^{-1}$.  In
\eqref{eq:mixed-unitary-spin}, compare choices $0$ and $d_i$ at
that place, using only these two choices at every other place.
Both resulting words belong to $W_{\rm CM}$, and their ratio is
$r_{i,p}^{d_i}$.  Hence
\begin{equation}\label{eq:mixed-unitary-root}
 r_{i,p}^{d_i}=1.
\end{equation}
This removes the case $d_i=1$.  For $d_i\ge3$, compare also the
choice $1$ with the choice $0$.  The ratio is $r_{i,p}$, and the
matched weight-one constituents show that these ratios agree on
shared simple ideals.  They are therefore supplied by the center
of the shared $\mathrm{SL}_{d_i}$ and can be removed there.
Such derived central corrections do not affect $W_{\rm CM}$.
The only remaining actions on the $V_i$ are therefore signs on
orthogonal dimension-four and unitary dimension-two blocks.
\end{proof}

\subsection{The graph and its auxiliary abelian varieties}

We now formulate the compatibility needed for these residual signs and
construct algebraic tensors that enforce it.  The graph and all its
auxiliary varieties depend only on the original Hodge representations,
before any element of the final group $C$ is chosen.

Form a finite graph $\Gamma$ as follows.  Its ordinary vertices are
the simple ideals of $\operatorname{Lie}(G_0)^{\rm der}_{\C}$
isomorphic to $\mathfrak{sl}_2$.  For every orthogonal factor with
$d_i=4$ and every place $p$ of $F_i$, put an edge $e=(i,p)$ between
the two ideals acting on that block.  For every unitary factor with
$d_i=2$ and every place $p$, put an edge between its simple ideal and
a new vertex $t_{i,p}$, used only for this edge.  We call these new
vertices \emph{torus vertices}.  Orthogonal dimension-three ideals
are included as ordinary vertices, even when no edge meets them.

An edge sign $\varepsilon_e$ will denote the scalar action on its
quadratic block after the preceding reduction.  To remove all these
signs using the shared derived group, we seek signs $s_v$ at the ordinary
vertices, and sign $1$ at each torus vertex, such that
\begin{equation}\label{eq:mixed-sign-target}
 \varepsilon_e=s_us_v
 \qquad(e\text{ with endpoints }u,v).
\end{equation}
Indeed, a central sign in $\mathrm{SL}_2$ acts on its standard module
by that sign.  The orthogonal dimension-four standard module is the
tensor product of the two endpoint modules, while the unitary
dimension-two standard module has just the ordinary endpoint's derived
action.  The same center acts trivially on the orthogonal
dimension-three standard module.

The necessary compatibility is that the edge signs multiply to one
around every cycle and along every path between torus vertices.
These relations are also sufficient: path products from a base vertex
give the signs $s_v$, as proved in Lemma~\ref{lem:mixed-node-signs}.
The role of the auxiliary tensors is to establish precisely these
relations for every normalized tensor automorphism.

We therefore use every simple cycle of $\Gamma$,
including a cycle formed by two parallel edges, and every simple path
with both endpoints torus vertices.  There are finitely many.  Fix one
such path or cycle $\gamma$.  All its edges have the same real field $F$ and
the same place $p$ under Lemma~\ref{lem:mixed-shared}.  For each
embedding $\tau:F\hookrightarrow\C$, choose a Galois transport
$\gamma_\tau$ of the path or cycle from $p$ to $\tau$.  It uses the
same list of original factor indices, although its simple ideals have
been transported.  At the original place choose the identity transport,
so that $\gamma_p=\gamma$.

Viewing each edge as an interval, define the \emph{grid} to be the
Cartesian product of these graphs,
\[
 \mathscr G_\gamma=
       \prod_{\tau:F\hookrightarrow\C}\gamma_\tau.
\]
Its vertices are tuples of vertices, one at each place; a grid step
changes one coordinate along an edge.  A \emph{box} is a product
\(\prod_\tau e_\tau\), where $e_\tau=(I(\tau),\tau)$ is one
chosen edge of $\gamma_\tau$.  Its vertices choose one endpoint of
each edge.  A square face of a box varies the endpoints at two distinct
places and freezes all other vertices.  At any torus endpoint we may
select either of its two character lines; include all such choices in
the finite construction.

\begin{lemma}[Realization of a box]\label{lem:mixed-box}
For each box there is an auxiliary abelian variety with weight-one
realization $W_{\Box}$ and an algebraic tensor
\[
 \mathcal P_{\Box}\in(\End W_{\Box})^{\otimes2}
\]
as in Proposition~\ref{prop:auxiliary-spin}.  At a chosen complex
embedding, $W_{\Box}$ contains the tensor representation attached
to every vertex of the box: the standard two-dimensional representation
at an ordinary vertex and the selected character line at a torus
vertex, tensored over the places of $F$.

Adjoining all these finitely many $W_{\Box}$ to the original system
does not change its Hodge Lie algebra.  Vertex representations with
the same choices, whether occurring in different boxes or in an
original $W_i$, are identified by algebraic weight-one correspondences.
\end{lemma}

\begin{proof}
Choose a totally real Galois closure $F'\subset\R$ of the embedded
field $F$.  Every embedding $q:F\hookrightarrow\C$ then takes values
in $F'$.  For each such $q$, take the four-dimensional quadratic block
of factor $I(q)$ and extend scalars to $F'$ through $q$.
In the unitary case this is the underlying
quadratic space of its two-dimensional Hermitian block.  Let $V'$ be
their orthogonal sum together with sufficiently many negative Tate
directions.

At every real embedding of $F'$, exactly one of the chosen blocks
has a period plane: composition with that embedding permutes the
places $q$, and their active places agree under the common-field
identification.  All other blocks are negative definite.  Thus $V'$
has signature $(2,d'-2)$ at every real embedding.  Choose the number
of negative directions so that $d'$ is sufficiently large and
$d'\equiv5\pmod8$.  The Witt index over $F'$ is at least two.
Indeed the nonarchimedean anisotropic dimension bound and the
Hasse principle first give a hyperbolic plane; the residual form
still has positive and negative directions at every real place and,
in sufficiently large dimension, gives a second hyperbolic plane.

Proposition~\ref{prop:auxiliary-spin} applies to
\[
 W_{\Box}=\operatorname{Res}_{F'/\Q}C^+_{F'}(V')
\]
and supplies its tensor $\mathcal P_{\Box}$.  The spin decomposition
at the chosen embedding gives every endpoint choice in the
four-dimensional blocks.  The extra negative Tate directions supply
only multiplicities, since the Hodge group acts trivially on them.
This proves the assertion about vertex representations.

Here is the rational group action underlying this construction.
Write $Q_i$ for the four-dimensional quadratic space over $F$ of
the edge's factor, and put $Q_{i,q}=Q_i\otimes_{F,q}F'$.  Extension
of scalars through $q$ gives a rational group homomorphism
\[
 \operatorname{Res}_{F/\Q}\mathrm{SO}(Q_i)
 \longrightarrow
 \operatorname{Res}_{F'/\Q}\mathrm{SO}(Q_{i,q}).
\]
Compose these maps with the factor actions of $G_0$, take their
orthogonal sum over $q$, and act trivially on the added Tate
directions.  This defines
\[
 \rho_{\Box}:G_0\longrightarrow
       \operatorname{Res}_{F'/\Q}\mathrm{SO}(V').
\]
The fiber product of $G_0$ with the corresponding spin covers,
using all the maps $\rho_{\Box}$, is a rational group mapping
onto $G_0$ with finite kernel.  The original Hodge circle together
with its auxiliary KS actions lies in this fiber product.
Consequently the enlarged Hodge group projects onto $G_0$ with
finite kernel and has the same Lie algebra.  Matching vertex
choices therefore have the same Lie representation in all their
occurrences.  Connectedness gives the same group representation,
and the weight-one correspondence argument of
Lemma~\ref{lem:mixed-constituents} makes their identifications and
multiplicity projections algebraic.
\end{proof}

Fix all the varieties in Lemma~\ref{lem:mixed-box}, and from now on
let $G,C$ denote the Hodge and algebraic tensor groups of this entire
finite collection.  We identify the Lie algebra of $G$ with that
of $G_0$.  In particular this collection satisfies the hypothesis
of the preceding normalization argument.  The preceding lemmas apply
to this final collection.
For a grid vertex $u$, let $R_u$ be the abstract tensor product of
its endpoint representations.  Identify it with a multiplicity copy
in every occurrence by weight-one Hodge embeddings.  Every
intertwiner between such copies is algebraic by
Lemma~\ref{lem:mixed-box}.  Hence any $c\in C(\C)$ induces the
same operator $c_u$ on $R_u$ through all these identifications.
We may therefore choose multiplicity copies in a box using its
Clifford tensor decomposition without changing $c_u$.

\subsection{The signs of a normalized tensor automorphism}

Now choose $c\in C(\C)$ and normalize it by
Lemma~\ref{lem:mixed-central}, using the CM part of this fixed final
collection.  For an edge $e=(i,p)$ of $\Gamma$,
let
\[
 \varepsilon_e\in\{1,-1\}
\]
be its scalar action on the associated quadratic block of $V_i$.
For a unitary edge this is a sign by
\eqref{eq:mixed-unitary-root} with $d_i=2$, and standard and dual
have the same sign.

At each four-dimensional block, a lift to
$\mathrm{Spin}_4=\mathrm{SL}_2\times\mathrm{SL}_2$ acts on the
two endpoint modules.  The normalized trivial adjoint action makes
each semisimple $\mathrm{SL}_2$ lift central, so its endpoint scalar
is a sign.  In the orthogonal case the two signs multiply to
$\varepsilon_e$.  In the unitary case one endpoint is a torus:
if $a_e\in\{1,-1\}$ is the ordinary endpoint sign and $\lambda_e$
is the selected torus character scalar, then
$a_e\lambda_e=\varepsilon_e$.  Thus $\lambda_e$ and its inverse,
the opposite character scalar, are signs as well.
These endpoint signs may at present
depend on the edge.  On a word in an original $W_i$, the action of
$c$ is the product of the endpoint signs selected at its places.
An overall sign of the spin lift can be absorbed into one block.

These are the edge signs in \eqref{eq:mixed-sign-target}.  We will
prove their cycle and torus-endpoint relations by transporting a sign
from one grid step to every parallel step, independently of the frozen
vertices.

\subsection{Transporting a sign across the grid}

We first isolate a linear-algebra observation that will justify the
subspaces used in the geometric tensor calculation.

\begin{lemma}[Supports of an invariant pairing]\label{lem:mixed-support}
Let $A,B$ be finite-dimensional complex vector spaces, let
$U\subset A$ and $V\subset B$, and let $t\in U\otimes V$ induce
an isomorphism $U^\vee\to V$.  If automorphisms $a$ of $A$ and
$b$ of $B$ satisfy $(a\otimes b)t=t$, then $a(U)=U$ and $b(V)=V$.
If moreover $a|_U=\varepsilon\,\mathrm{id}$ for
$\varepsilon\in\{1,-1\}$, then
$b|_V=\varepsilon\,\mathrm{id}$.
\end{lemma}

\begin{proof}
The images of the contraction maps $B^\vee\to A$ and
$A^\vee\to B$ defined by $t$ are exactly $U$ and $V$.
Thus preservation of $t$ preserves both subspaces.  The last
assertion follows from the perfect pairing, since
$\varepsilon^{-1}=\varepsilon$.
\end{proof}

Fix one of the paths or cycles used to construct a grid.  For a step
along its edge $e=(i,p)$ at place $p$, freeze the vertices at every
other place.  If the local endpoint spaces are $X,Y$ and the tensor
product of the frozen spaces is $T$, the two grid representations are
$X\otimes T$ and $Y\otimes T$.  Inside the Hom space between them,
consider
\begin{equation}\label{eq:mixed-local-hom}
 \Hom(X,Y)\otimes\mathrm{id}_T.
\end{equation}
The reverse direction uses $\Hom(Y,X)\otimes\mathrm{id}_T$.
Neither subspace is being assumed algebraic on its own.

\begin{lemma}[Propagation across a box]\label{lem:mixed-propagation}
For the normalized element $c$, the action on
\eqref{eq:mixed-local-hom}, in either direction, is the scalar
$\varepsilon_e$, independently of the frozen vertices at the
other places of the grid.
\end{lemma}

\begin{proof}
First choose every frozen vertex at an endpoint of the transported
edge of factor $i$.  The corresponding tensor words occur in the
original $W_i$.  Its endpoint-sign description shows that the action
on \eqref{eq:mixed-local-hom} is $\varepsilon_e$ in both directions:
the frozen signs cancel, and the ratio of the two endpoint signs is
their product.

It remains to move one frozen vertex across an edge at a different
place $q$.  Choose a box containing these two steps, with the other
places fixed.  Denote its four vertex representations by
\[
 R_{00}=X\otimes Z\otimes T_0,\quad
 R_{10}=Y\otimes Z\otimes T_0,\quad
 R_{01}=X\otimes Z'\otimes T_0,\quad
 R_{11}=Y\otimes Z'\otimes T_0.
\]
Project $\mathcal P_{\Box}$ to
\begin{equation}\label{eq:mixed-cross-pairing}
 \Hom(R_{00},R_{10})\otimes\Hom(R_{11},R_{01}).
\end{equation}
The second $p$-direction has been reversed.  This projection is
algebraic, since its vertex projections are weight-one Hodge
endomorphisms.  Figure~\ref{fig:mixed-grid} records the two directions.

\begin{figure}[htbp]
\centering
\setlength{\unitlength}{1pt}
\begin{picture}(310,110)
 \put(30,85){\makebox(0,0){$R_{00}$}}
 \put(260,85){\makebox(0,0){$R_{10}$}}
 \put(30,20){\makebox(0,0){$R_{01}$}}
 \put(260,20){\makebox(0,0){$R_{11}$}}
 \put(52,85){\vector(1,0){186}}
 \put(238,20){\vector(-1,0){186}}
 \put(30,66){\vector(0,-1){28}}
 \put(260,66){\vector(0,-1){28}}
 \put(145,99){\makebox(0,0){$p$-step: $\varepsilon_e$}}
 \put(145,4){\makebox(0,0){reverse $p$-step: $\varepsilon_e$}}
 \put(16,52){\makebox(0,0)[r]{$q$}}
 \put(274,52){\makebox(0,0)[l]{$q$}}
 \put(145,52){\makebox(0,0){$\mathcal P_{\Box}$ pairs the local Hom subspaces}}
\end{picture}
\caption{A square face of a box in the grid of transported paths or
cycles.  Only the distinct place coordinates $p$ and $q$ vary.
The projected tensor pairs the local Hom subspace represented by the
upper $p$-arrow with the one represented by the reversed lower arrow.
Preservation of this perfect pairing transports the scalar sign when
the frozen $q$-vertex changes.  The vertical arrows indicate that change
of vertex, not a chosen linear map.}
\label{fig:mixed-grid}
\end{figure}

We verify that the tensor in \eqref{eq:mixed-cross-pairing} is a
perfect pairing between the two local Hom subspaces.  Only the
$p$-block of $V'$ contributes.  Clifford multiplication by a vector
in that four-dimensional block interchanges its two half-spin
spaces and acts as the identity on spectator slots, up to fixed
grading signs.  In a graded Clifford tensor model these signs are
the chirality operators on preceding even-dimensional slots; after
choosing the spectators they are scalars.  The volume element in
the total odd-dimensional Clifford module likewise acts as a
scalar.  Choose multiplicity copies in this Clifford tensor
decomposition and transport the projected tensor to the fixed
tensor models $R_u$.  Its spectator factors are then identities;
the transported actions are the already fixed operators $c_u$.

For the two two-dimensional half-spin spaces $X,Y$, the vector
representation is $X\otimes Y$, and Clifford multiplication gives
all of $\Hom(X,Y)$ and all of $\Hom(Y,X)$.  More explicitly, with
the skew-form identifications $j_X:X\to X^\vee$ and
$j_Y:Y\to Y^\vee$, the two maps belonging to $x\otimes y$ are,
up to nonzero scalars,
\[
 y\otimes j_X(x),\qquad x\otimes j_Y(y).
\]
The inverse quadratic metric therefore gives a nondegenerate tensor
between the two Hom spaces, with identities on the frozen factors.
If an endpoint is a torus character line in $Y$, project to that
line.  The first Hom space consists of maps from $X$ into the selected
line, and the second of maps from that same line to $X$.  In the
displayed formulas these use respectively the selected character
component of $y$ and its opposite under $j_Y$.  They remain perfectly
paired by the metric.  Thus the projected tensor has
exactly the local Hom subspaces as its minimal tensor supports.

Since $c$ preserves $\mathcal P_{\Box}$ and all the vertex
projections, Lemma~\ref{lem:mixed-support} first makes these Hom
subspaces invariant and then transfers the sign $\varepsilon_e$
from the upper one to the lower one.  Reversing the roles of the
two directions gives the reverse propagation as well.  Moving the
frozen places successively along their paths proves the assertion
for every vertex choice.
\end{proof}

\begin{lemma}[Cycle and endpoint relations]\label{lem:mixed-cycle}
The edge signs of $\Gamma$ satisfy
\begin{equation}\label{eq:mixed-cycle-signs}
 \prod_{e\in\gamma}\varepsilon_e=1
\end{equation}
for every cycle $\gamma$ and for every path $\gamma$ whose two
endpoints are torus vertices.
\end{lemma}

\begin{proof}
It suffices to use simple cycles and simple paths, since cancellation
of repeated edges reduces the general case to them.  Fix one such
path or cycle and a place $p$.  Freeze all other places at the
starts of their transported paths.  At the initial grid vertex,
the action of $c$ is scalar: all these endpoint choices occur in
the original representation belonging to the first edge.

Follow the path at place $p$.  Suppose $c$ is scalar $a$ on its
current vertex representation.  By
Lemma~\ref{lem:mixed-propagation}, conjugation on the local Hom
space for the next edge $e$ is $\varepsilon_e$.  Hence
\[
 c_{\rm next}\,f=\varepsilon_e a f
 \qquad
 \text{for every }f\in\Hom(X,Y)\otimes\mathrm{id}_T.
\]
The images of these maps span the next vertex representation.
Thus its scalar is $\varepsilon_e a$.  Around a cycle one returns
to the same representation and obtains
\eqref{eq:mixed-cycle-signs}.  For a torus-to-torus path, both
endpoint words consist only of character slots, including at the
frozen places.  They belong to $W_{\rm CM}$, so their scalars are
both one.  The same calculation proves the endpoint relation.
\end{proof}

\begin{lemma}[Vertex signs]\label{lem:mixed-node-signs}
There are signs $s_v$ at the vertices of $\Gamma$, with $s_v=1$
at every torus vertex, such that
\[
 \varepsilon_e=s_us_v
 \quad\text{for every edge }e\text{ with endpoints }u,v.
\]
\end{lemma}

\begin{proof}
In each connected component choose a base vertex, choosing a torus
vertex when one is present, and give it sign one.  Define the sign
of any other vertex as the product of the edge signs along a path
from the base.  The cycle relations in
Lemma~\ref{lem:mixed-cycle} make this independent of the path.
If another torus vertex occurs, the endpoint-path relation gives
it sign one.  The edge formula follows immediately.  Isolated
vertices may also be assigned sign one.
\end{proof}

\subsection{Completion of the mixed-product argument}

\begin{proof}[Proof of Theorem~\ref{thm:main}]
Take $c\in C(\C)$ for the final collection fixed after
Lemma~\ref{lem:mixed-box}.  Lemma~\ref{lem:mixed-central} removes
an element of $G(\C)$ so that the only remaining actions on the
$V_i$ are the edge signs of $\Gamma$.  By
Lemma~\ref{lem:mixed-node-signs}, these are products of signs
assigned to the shared $\mathfrak{sl}_2$ ideals, with sign one at
the torus vertices.  Use those signs in the centers of the
corresponding simply connected derived factors.  On an orthogonal
dimension-four block their action is the product of the two
endpoint signs; on a unitary dimension-two standard block it is
the ordinary-vertex sign.  Their action on orthogonal
dimension-three blocks is trivial.  They therefore remove every
remaining action on the $V_i$.  This proves
\eqref{eq:mixed-group-goal} over $\C$, and hence over $\Q$.

Every rational Hodge tensor in the $V_i$ is now fixed by $C$ and
is therefore algebraic.  Apply the algebraic K\"unneth projectors
to each factor of $S_1\times\cdots\times S_m$.  Besides the
transcendental summands in degree two, the pieces are generated by
the unit, divisors, and the point class, all algebraic.  Thus every
Hodge class in a K\"unneth summand is a sum of the algebraic
transcendental tensors just obtained tensored with these elementary
classes.  Products and the projectors realize them in the required
codimension.  This proves surjectivity of the rational cycle-class
map for every ordered list, including repetitions, and every
codimension in the statement.
\end{proof}
